\documentclass[10pt]{article}
\usepackage{amsmath,amsfonts,amssymb,mathtools,mathrsfs,bm,bbm,array,xcolor,geometry}
\allowdisplaybreaks
\setlength{\parindent}{0pt}
\setlength{\parskip}{4pt}
\emergencystretch=4em

\geometry{paperwidth=215mm,paperheight=290mm,hmargin=15mm,vmargin=16mm}
\begin{document}
\section*{P0008 Daily Transport and Processing of Diagnostic Samples and Results}
\subsection*{Mathematical formulation}
\begin{equation}
\min\quad \frac1T\sum_{t=1}^T\sum_{i\in I}\sum_{k=1}^7s_{ik}(t)+\sum_{i\in I}\sum_{k=1}^7\rho_k s_{ik}(T).
\end{equation}

\begin{align}
d(t)&=\sum_{(i,j)\in A}D_{ij}x_{ij}(t) && t\in\mathcal T,\\
\sum_{t\in\mathcal T}d(t)&\leq B,&\sum_{t\in\mathcal W_1}d(t)&\leq B_1,\\
\sum_{j\ne1}x_{1j}(t)&=\sum_{j\ne1}x_{j1}(t)\leq c_t &&t\in\mathcal T,\\
\sum_{j\ne i}x_{ij}(t)&=\sum_{j\ne i}x_{ji}(t)=y_i(t) && i\ne1,t\in\mathcal T,\\
x_{1j}(t)&\leq y_1(t),\quad y_1(t)\leq\sum_{j\ne1}x_{1j}(t)&&j\ne1,t\in\mathcal T,\\
x_{jL}(t)&=x_{Lj}(t)=0&&j=2,\ldots,n,t\in\mathcal T,\\
0\leq h_i(t)&\leq H_t-1 && i\in V_t\setminus\{1\},t\in\mathcal T,\\
h_j(t)-h_i(t)&\geq1-H_t(1-x_{ij}(t))&& i,j\in V_t\setminus\{1\},i\ne j,t\in\mathcal T,\\
D_{i1}\leq g_i(t)&\leq R_t-D_{1i}&& i\in V_t\setminus\{1\},t\in\mathcal T,\\
g_j(t)-g_i(t)&\geq D_{ij}-R_t(1-x_{ij}(t))&&i,j\in V_t\setminus\{1\},i\ne j,t\in\mathcal T.
\end{align}

\begin{align}
0\leq f_{ik}(t)&\leq M z_{ik}(t),\\
f_{ik}(t)&\leq s_{i,k-1}(t-1),\\
f_{ik}(t)&\geq s_{i,k-1}(t-1)-M(1-z_{ik}(t)).
\end{align}

\begin{align}
s_{ik}(0)&=\bar s_{ik} && i\in I,k=1,\ldots,7,\\
f_{i1}(t)&=u_i(t) &&i\in I,t\in\mathcal T,\\
f_{ik}(t)&=b_{ik}(t)+\sum_{\tau=1}^{t}\alpha_k(\tau,t)f_{i,k-1}(\tau)&&i\in I,k\in\{3,5,7\},t\in\mathcal T,\\
s_{ik}(t)&=s_{ik}(t-1)+f_{ik}(t)-f_{i,k+1}(t)&&i\in I,k=1,\ldots,7,t\in\mathcal T,\\
s_{ik}(t),f_{ik}(t),d(t)&\geq0,\quad x_{ij}(t),y_i(t)\in\{0,1\}&&\text{all defined indices}.
\end{align}

\subsection*{Definitions and model scope}
A health-service manager plans courier movements over a finite horizon of $T$ days to send diagnostic samples to a laboratory and return results to their originating facilities. The days $\mathcal T$ run from one through $T$, and locations $V$ run from one through $n+1$: location one is the district hub, locations two through $n$ are other health facilities, and the last location $L$ is the molecular laboratory. The origins $I$, numbered one through $n$, include the hub itself; its samples also undergo transport, processing and waiting-cost accounting. Directed connections $A$ join every pair of distinct locations. Distances $D_{ij}$ are nonnegative and symmetric, self-distances are zero, and a direct distance cannot exceed the two-leg distance through the hub. Each day's courier count $c_t$, maximum non-hub stops per route $H_t$ and maximum route mileage $R_t$ are fixed. A location whose round trip from the hub exceeds that day's route-mileage limit cannot be served that day.

The manager uses binary states $x_{ij}(t)$ to choose daily connections and $y_i(t)$ to record visits, with one meaning used or visited and zero meaning otherwise. Departures from the hub equal returns and cannot exceed the day's courier count. At each non-hub location, both the number of arrivals and the number of departures equal its visit state, so it can be visited at most once per day. The hub is marked visited exactly when at least one route departs; it cannot be marked visited without a route or left unmarked when a route departs. Travel in either direction between the laboratory and non-hub health facilities is forbidden, so laboratory trips run separately through the hub. Nonnegative daily mileage $d(t)$ equals the sum of the distances of chosen connections. Total mileage cannot exceed $B$ over the full horizon or $B_1$ over the first-week days $\mathcal W_1$, which cover the shorter of the first seven days and the horizon.

Each non-hub location within that day's service radius has a continuous order mark $h_i(t)$ between zero and the maximum stop count minus one. For any two distinct such locations, the destination mark minus the origin mark must be at least one on a chosen connection and at least one minus the maximum stop count on an unchosen connection. This limits stops on each route and prevents cycles disconnected from the hub. A continuous distance mark $g_i(t)$ lies between the distance from that location to the hub and the daily route-mileage limit minus the distance from the hub to that location. A chosen non-hub connection requires the destination mark to exceed the origin mark by at least its distance; for an unchosen connection, that required increase is reduced by the daily route-mileage limit. These ranges and pairwise checks remain in force for unvisited locations within the service radius, including connections forbidden by the laboratory rule. Locations beyond the daily radius have no order or distance marks that day, and all their visits and incoming and outgoing connections are zero. Their samples and results still remain in the stage-flow and waiting-cost accounts; any transport requiring their own visit moves nothing that day.

Samples and results pass through eight stages: awaiting collection at a facility, awaiting registration at the hub, registered at the hub, awaiting testing at the laboratory, results available at the laboratory, results awaiting registration at the hub, registered results at the hub, and results returned to the facility. Stage eight incurs no further waiting. For origin $i$ and stage $k$, end-of-day stock $s_{ik}(t)$ and the amount entering that stage $f_{ik}(t)$ may be nonnegative continuous quantities. Stocks are tracked for the first seven stages and entries for all eight. Initial stocks $\bar s_{ik}$ are nonnegative, and daily entry into stage one equals the known arrivals $u_i(t)$. At each of the first seven stages, closing stock equals previous-day stock plus today's entries minus today's entries into the next stage and cannot be negative.

Entries into stages two, four, six and eight require transport. The visit trigger $z_{ik}(t)$ equals the origin's visit state for stages two and eight and the laboratory's visit state for stages four and six. A triggered visit moves exactly all previous-day stock at the preceding stage; without a visit it moves nothing, so stock newly made ready that day is not moved along with it. A common quantity bound $M$ is at least the largest, over origins, sum of initial stocks across seven stages and arrivals across the horizon. Every transported amount is nonnegative, at most this bound times the visit trigger, and at most previous-day preceding-stage stock; it is also at least that stock minus this bound times the no-visit state.

Stages three, five and seven are entered upon processing completion. The fraction $\alpha_k(\tau,t)$ specifies how much of a batch entering the preceding stage on day $\tau$ completes on day $t$. It lies between zero and one, is zero before entry, totals at most one over the horizon for each batch, and may permit completion on the entry day. Today's entry into the processed stage equals the sum of preceding-stage entries on days one through today times their corresponding completion fractions, plus the scheduled completion $b_{ik}(t)$ from stock already being processed at the start. This nonnegative quantity follows the same processing calendar for pre-horizon batches and totals no more than the origin's initial preceding-stage stock. It is zero for an initially empty system but cannot be ignored for a nonempty one. The manager seeks the lowest total of closing stocks for all origins across the first seven stages and all days, divided by the horizon length, plus terminal stocks weighted by their nonnegative stage penalties $\rho_k$. Arrivals, processing calendars, initial stocks and scheduled historical completions are fixed in advance; closing stocks result from the chosen transport and the stage flows described above. Visits have no additional frequency requirement, routes are not restricted to a preselected menu, and this finite-horizon arrangement does not represent the best policy for indefinite stochastic operation.

\subsection*{Inline notation and relations}
$\mathcal T=\{1,\ldots,T\}$; $V=\{1,\ldots,n+1\}$; $2,\ldots,n$; $L=n+1$; $I=\{1,\ldots,n\}$; $A=\{(i,j):i,j\in V,i\ne j\}$; $D_{ij}=D_{ji}\geq0$; $D_{ii}=0$; $c_t,H_t,R_t$; $B,B_1$; $\mathcal W_1=\{1,\ldots,\min(7,T)\}$; $V_t=\{1\}\cup\{i\in V\setminus\{1\}:2D_{1i}\leq R_t\}$; $V\setminus V_t$; $V_t\setminus\{1\}$; $I$; $H_t\ge1$; $c_t\ge0$; $R_t$; $D_{ij}\leq D_{i1}+D_{1j}$; $u_i(t)$; $\bar s_{ik}\geq0$; $k=1,\ldots,7$; $\rho_k\geq0$; $k\in\{3,5,7\}$; $\alpha_k(\tau,t)\in[0,1]$; $\tau$; $t$; $\alpha_k(\tau,t)=0$; $t<\tau$; $\alpha_k(t,t)>0$; $b_{ik}(t)\geq0$; $\sum_t b_{ik}(t)\leq\bar s_{i,k-1}$; $\tau\leq0$; $M\geq\max_i[\sum_{k=1}^7\bar s_{ik}+\sum_tu_i(t)]$; $x_{ij}(t)$; $y_i(t)$; $h_i(t),g_i(t)$; $d(t)$; $s_{ik}(t)$; $f_{ik}(t)$; $k=1,\ldots,8$; $y_i=\max_jx_{ij}$; $i=1$; $z_{ik}(t)=y_i(t)$; $k\in\{2,8\}$; $z_{ik}(t)=y_L(t)$; $k\in\{4,6\}$; $i\in I,t\in\mathcal T,k\in\{2,4,6,8\}$; $f_{ik}(t)=z_{ik}(t)s_{i,k-1}(t-1)$.
\end{document}